\chapter*{前\quad 言}
\addcontentsline{toc}{chapter}{前言}
\markboth{前言}{前言}

AI可以在一场会议结束之前，写出十份足以发起下一场会议的方案。方案越来越便宜，会议、施工、试验和承担后果的人，却没有按同样的速度增加。关于人工智能的争论，往往跳过了这段距离。看好技术的人说，模型还会更强，成本还会更低；怀疑技术的人说，落地总比演示慢。两边都能找到证据，却共用了一个前提：只要智能足够多、足够便宜，其余问题就会依次让路。本书从这个前提开始。

生成已经变便宜了。过去几年，模型、数据、训练方法与推理计算共同推动了能力增长，写作、图像和代码生成进入日常工作。但降下来的成本没有均匀扩散到整条链。一份结构完整的方案几十秒就能写出，一条生产线仍要花几个月到几年才能运行：设备要采购，人员要培训，流程要审批，责任要重新划分。瓶颈从生成端移到了别处。把这次移动读成整个世界都在加速，判断便从这里开始偏。

第一个要检查的前提是智能可以被单独衡量，并且沿着一条单一的轴持续变强。这个前提在工程上有来历：模型规模、数据量、算力确实在一段时间里和性能同步上升，以至于外推变成了一种习惯。但智能不是一个标量。识别、推理、记忆、规划、动手、共情、承担，这些能力互不替代，各自要求的条件也不同；机器在其中的一些上已经超过人，在另一些上还很差，而这种差异不是暂时的。一个只在文本上训练出来的系统，它的能力边界不会因为文本变多而自动扩到物理操作、长期承诺和后果承担上去。把智能当成单一刻度，就会得出"再强一点就全都解决了"的推论，而这个推论至今没有被兑现过。

第二个前提是复制的低成本可以推广到执行。文档、图像和程序可以被反复复制，数字化也确实替代了地图、胶片、唱片和柜台上的大量操作。但图纸的复制与桥梁的建造消耗不同的资源。一座桥的图纸可以分发一万份，每一座桥仍要单独施工；混凝土的强度发展取决于配合比、温湿度和养护条件，工程中常用的28天龄期是检验基准；药物候选可以快速生成，疗效与长期安全性仍要经过试验和观察。设计、筛选和验证工具能够改变这些过程的效率，过程所要求的材料、暴露时间和现实反馈仍须逐项满足。

第三个前提是更优的方案会自动胜出。工程史里充满反例。更好的技术路线输给更成熟的供应链、更庞大的既有安装量和更低的切换成本，是常态而不是例外。一个方案要落地，必须穿过资源、时间和既有结构这三重约束，而优化只作用在其中很小的一段上。与它相邻的还有一组关于人的误解。流行的说法是替代：机器会抢走工作，或者机器会把人解放出来。实际发生的事情更细碎，也更难对付。被置换的往往不是人，而是人与成果之间的那套证明关系。一个人会写、会算、会总结，这些能力使他被需要，也使他的位置不言自明；当同样的产出可以由系统在几十秒内生成，能力还在，证明力却没了。这个问题在第五、六两章分别从塌陷机制和价值迁移两个方向处理。

这几个前提之所以顽固，是因为它们都能在同一个层面上自洽。在信息层面推理，复制确实免费，优化确实有效，能力确实可以外推；每一条在自己的语境里都成立，推广出去才出错。要判断它们在哪里失效，必须往下走一层，问一些更基础的问题：智能是什么，它从哪里来，它为什么会长成现在这个样子，它进入物理世界和生命世界时会撞上什么。这些问题在物理学、生物学和信息科学里都有现成的答案，只是平时不在同一个语境里出现。本书做的事情，是把这些答案接在一起，用它们来检查上面那几个前提。

接这些答案用的中心概念是黏滞。黏滞不是慢。慢是一个速度判断，黏滞是一个代价判断。一个系统从状态 A 转到状态 B，中间要穿过的东西才是黏滞；它问的不是这一步花了多久，而是这一步消耗了什么，留下了什么，出错之后能不能退回来。黏滞也不是保守或者技术恐惧。数字支付、导航、工业仿真、自动化实验和代码补全，都确实降低了某些环节的黏滞，本书不否认这一点。真实情况更分散：不同环节的黏滞程度差别很大，同一项技术在不同环节的沉降速度也不同。需要反复问的是同一件事，这项技术降低了哪一类黏滞，又把黏滞转移到了哪里。

任何一次现实状态转移，至少要检查三类约束。推动改变所需的材料、能源、设备和资本构成资源势垒（$R$）；反应、生长、测试和协调占用的过程时间构成时间刚性（$T$）；改变沿依赖网络传导时遇到的同步、隔离和回滚困难构成拓扑耦合（$G$）。三者共同描述状态转移所处的位置。工程改进既可能以增加资源换取时间，也可能通过新工艺同时降低多项成本；局部降阻还可能把负担转给维护、验证或其他承担者。判断技术效果，需要同时核对总成本的变化和成本的分布。

同一组约束在不同介质里的表现差别很大，所以本书把智能进入现实时面对的世界分成三种介质。\term{比特}是信息介质，复制、传播、修改和回滚的成本都低，它最先被生成式系统改造，也最容易让人产生万物皆可重构的印象；但它的低黏滞集中在生成层面，不在可信层面，生成不再成为问题之后，验证成为问题。\term{原子}是物质介质，制造、建筑、交通、能源、物流、芯片和基础设施都要面对质量、能量、摩擦、惯性、材料强度和空间约束；一个语言错误可以删除，一个机器人错误可能打碎物品、伤害人体或者损坏自身。\term{生命}是第三种介质，它不响应指令，只响应干预：一个改变进入生命系统之后会被识别、代偿、适应，甚至反弹；DNA 可以写成序列，基因功能却不等于文本意义。三种介质的区别是吸收机制的区别，比特把智能转译为扩散，原子把智能转译为做功，生命把智能转译为反馈。用同一种加速逻辑去驱动三者，必然出错。

全书七章，加一个附录。第一章处理时间感的落差：为什么生成的速度会被误当成做事的速度，以及黏滞为什么同时是阻力与稳定器。第二章给出物理疆界，智能从哪里来，物质重构要支付什么样的资源势垒、时间刚性和拓扑耦合代价；中心结论是算力可以优化路径，不能取消势差，可以压缩搜索，不能压缩过程，可以生成方案，不能替代状态转移。第三章把这些约束放进接口，看三种介质各自如何吸收智能输入，现实系统在过载时会绕行、堆积还是失稳，生成物又在什么条件下才算被确认。

第四章转向生命与身体，并把这套约束一直带到位移系统上。前三节处理生命系统为什么不响应指令而响应干预、机器身体为什么不能简单继承生物身体的能力；第四节说明智能为什么会分化成探索、控制、验证三种模态，而不是收敛成单一全能系统；第五节把这个分化落回到人身上，从任务能耗、反馈稳定性和行动准入三个方面讨论人的位置，区分当前的工程优势与需要制度确立的判断资格；第六、七两节把同一套账本用到地面交通和低空系统上，看代价如何在动力平台、既有基础设施与治理层之间转移；第八节给出这套分析的用法：管理多速率耦合的系统，目标不是求极值，是维持一个可行的流态区间。

第五章是全书重心：什么样的东西算存在，什么样的判断可以被依赖，一个表示什么时候可以跨过转换边进入行动，行动之后的后果回到哪里，什么样的边界才有资格被称为承受者。这一章的前两节把这些判据立起来，后四节顺着它们往下走，给出功能性、真实性与主体性三层依次发生的塌陷，以及它们如何叠成一个下沉的螺旋。从这里开始，问题不再只是什么能被生成，而是什么能被承担。第六章处理价值迁移：技术每一次降低旧黏滞，人的价值证明方式都会迁移一次，文字、印刷术、钟表、蒸汽机、计算机，每一次都是；生成成本下降以后，验证、照护、长期承诺与后果承担可能成为新的稀缺环节；这些环节能否获得报酬和社会承认，取决于分配制度。这一章的后半段把迁移的方向收成四条原则（阻尼本体论、不完备性原则、共生原则与递归优化原则），它们分别处理代价如何证明边界的存在、价值被压成目标函数时会损失什么、分工为什么是互补而不是竞争、调整为什么必须逐层进行，并说明它们为什么更接近一部宪法的条款，而不是一份技术方案。第七章收束到实践：个体如何保留现实粗糙度，组织如何重建信任和验证，文明为什么仍然需要高阻力前沿，以及自复制系统为什么会成为整套理论的压力测试。附录给出广义黏滞模型的完整形式、它在已有理论中的位置，以及把它变成可检验命题的路线。

这几章的次序不是任意的。前三章建立约束，第四章看约束如何在生命与身体上兑现成具体的代价，第五、六章处理约束被绕开之后留下的空缺，第七章回到行动。每一章都从上一章停下的地方接续，因此全书适合按顺序读；单独拆开读某一章也可以，但每一章的结论都以前面已经建立的那部分为前提，跳读时需要回到相应的前提。

本书以约束分析连接工程问题与价值问题。无效等待和冗余应当消除，执行与验证的成本需要降低，安全、信任和责任的功能则需要保留或由其他可靠机制接替。物理模型用于辨认相应的作用关系，具体部署的判断还要加入任务数据、组织条件与承担后果的人。

如果AI能替代越来越多的产出功能，人凭什么继续有位置。这个问题涉及能力，也涉及社会用什么尺度评价人。人的处境包含身体、时间、关系、判断和后果承担，产出只能描述其中一部分。承认这些内容的价值，需要相应的权利、分配和责任安排。自由不是没有黏滞，而是有能力选择、分配和承担黏滞。
